\topic{Mathematical destination: the facts behind the experiments}
\labelled{Translation preserves a heading.}{Glue each side of a regular octagon to the opposite side by a translation: corresponding points are the same point of the surface. A traveler exits at one point and resumes at its mate, keeping the same direction. There is \textbf{no reflection and no turn}. Position \emph{on the surface} and heading must both return for a closed trip. At a vertex, stop: the directional flow has no unique continuation there. [S1]}
\labelled{Horizontal trips form two families.}{Use an octagon with side length 2, centered at the origin, with horizontal and vertical sides. Put \(a=1+\sqrt2\), \(b=2+\sqrt2\). For \(-1<y<1\), rightward trips close after one edge crossing and have length \(2a\). For \(1<y<a\), a trip visits an upper and a lower strip, closes after two crossings, and has length \(2b\). The lower strip belongs to the same family. Each whole loop has the same length within its family, although its separate drawn segments vary. Lines through vertices are excluded. Children can compare with tracing strips, without square roots.}
\labelled{Eight corners become one exceptional point.}{Endpoint matching identifies all eight vertices. Each contributes \(135^\circ\); together they give \(1080^\circ\), or \emph{three full turns}. Away from this point the surface is Euclidean-flat, including across joined edges. Its one-face, four-edge, one-vertex description has Euler characteristic \(-2\), so the connected closed orientable surface has genus 2. The two-handle conclusion uses a classification theorem; it is adult context, not something proved by a sketch. [S1]}
\labelled{The three-square L has the same topology, a different geometry.}{Label its lower-left, lower-right and upper-left squares A, B and C. Right-edge crossing swaps A/B and fixes C; top-edge crossing swaps A/C and fixes B. The twelve square-corner sectors form one \(1080^\circ\) cone point. Again the genus is 2, but this is a different flat metric. The L is not a regular octagon cut into three equal squares. [S2, S3]}
\labelled{Rational directions close in the square model if they avoid corners.}{A rational direction repeats its local coordinates on the unit square after a fixed displacement block. That block permutes the finitely many square labels, so some repetition returns the entire state. Returning to the same local dot in a different square is insufficient. The proof is for the square-tiled model; it does \textbf{not} supply a rational/irrational slope test for the regular octagon. Finite drawings prove neither density nor ergodicity.}
\labelled{Two different octagon pictures must stay separate.}{Translated regular Euclidean octagons can overlap in a route drawing; they do not tile the ordinary plane. The usual hyperbolic genus-two octagon tiling uses a different metric, with \(45^\circ\) corners when eight meet smoothly. Our octagon has \(135^\circ\) corners and a cone point. [S4]}
\topic{What children are meant to discover}
Make and check a portal crossing; find and compare complete trips; notice that moving the start can change a trip or preserve its total length; chase corner identifications; then use an exact square lab to explain return through a finite set of states. A drawing is an \textbf{experiment}; a proposed rule is a \textbf{conjecture}; a matching argument or exact calculation is an \textbf{explanation}.
\par\begin{tabularx}{\linewidth}{@{}p{1.06in}X@{}}
\toprule
Entry & Readiness, independent of grade \\
\midrule
Grades 3--5 core & Match edge letters and arrows; move a ruler without turning; compare positions along equal edges; retain a start dot and travel arrow. Adult reading/recording is fine. No degree measure, fractions, algebra or square roots is required. \\
Later routes & Corner work needs endpoint matching and sector counting. The L lab needs square labels, grid positions and a repeated direction step. The general return proof needs reversible moves; adult formulas are not prerequisites. \\
\bottomrule
\end{tabularx}
\newpage
\topic{Prepare a small, usable kit}
\labelled{Per pair or trio.}{Student pages for the chosen route, 3--5 sheets of tracing paper or translucent baking paper, a ruler, pencil/eraser, two colored pencils, plain paper, one small movable counter, and a separate short paper direction arrow. Have scissors for the corner sectors; an adult can cut. A spare printed copy preserves an example while the next trial is made.}
\labelled{Printing and space.}{Print single-sided on US Letter: pp. 1--2 per pair for a first visit; p. 3 and tracing sheets for developing; pp. 4--5 for corners, keeping p. 4 uncut; pp. 6--9 only as chosen. Match traced copies to their original size; check the p. 5 scale before cutting. The optional three-copy p. 3 layout needs about 12 by 8 inches of desk. Never stretch the octagon's axes independently. No child needs to fold it into a two-handled object.}
\labelled{Staffing and roles.}{Use pairs wherever possible. The traveler chooses or draws the path; the checker points to the matching edge, locates the corresponding point, and checks the unchanged direction. Swap after a complete trial. In a trio, rotate a recorder who preserves one path and its start arrow. An adult may read, steady tracing paper, or record a dictated observation; children should still choose and check the path.}
\labelled{Before children arrive.}{Select a short route and read its answers. Trace one working octagon, including edge labels, matching arrows and a clear top mark. Keep a loose direction arrow next to it. Separate unused pages so that a thick packet does not become a completion demand. Preparation needs ordinary printing and tracing, not a custom three-dimensional model.}
\labelled{Physical rehearsal still needed.}{Try a slanted-edge transfer with the actual print, ruler and tracing paper. Check that labels remain visible through the paper; corresponding fractions of an edge are easy to match; children can keep the direction arrow fixed while moving a counter; corner marks and grid dots remain distinguishable. Test sector cutting and handling. \textbf{These materials are unpiloted. Digital checks do not establish physical fit or classroom usability.}}
\topic{A shared launch: slide the map, keep the heading}
\labelled{Handle first, about 2 minutes.}{Let children move a counter and trace a little of the octagon. Ask one child to find matching edges. Point out that the small printed edge arrows tell which points correspond; the loose arrow tells where the traveler is going. They do different jobs.}
\labelled{Show one legal crossing, about 3 minutes.}{Use the student packet's non-task visual. On a spare octagon, approach a nonvertex point on the right vertical edge along a short slanting line. Stop at the edge. Match its position on the left edge, measured from the same end of the identically directed arrows. Move the counter to its mate and draw the next short segment \emph{parallel to the first}, traveling in the same direction. Do not finish a loop or announce the length rule.}
\labelled{Show why the pieces are straight.}{Trace the octagon on transparent paper. Slide, do not turn or flip, the copy until its left edge lies on the original right edge. The arrival segment and continued segment now lie on one ruler line. Mark the copy 2. In the original drawing, the resumed piece belongs at the matched left-edge point. Physical overlap never tells us which surface points are equal; edge labels do.}
\labelled{Children demonstrate the rule.}{A volunteer makes one crossing at a different point. The checker touches both corresponding points, then compares travel arrows. Say: Same place on the matching arrow, same heading. If you reach a corner, stop. A trip is back only at its own starting point with its starting heading. Give pairs their first selected task.}
\labelled{Keep one visible record.}{A path with its starting dot, direction arrow and numbered crossings is enough. Do not demand a crossing tally, code, distance sum and oral explanation simultaneously. Save the paper so another pair can check it.}
\newpage
